\documentclass[11pt]{article}
\usepackage{icmlko}
\begin{document}

\title{짝지은 효과를 짝 없는 자로 설명하려 했다 --- 노트 492 $\cdot$ 493 취소}
\author{월드모델 실험실}
\date{노트 495 · 2026년 8월 3일}
\maketitle

\begin{abstract}
노트 494 가 덮음 이동 스물다섯 짝을 세면서 \emph{행 수를 안 봤다.}
되뽑기 구간을 붙이니 \textbf{$24$ 개 중 잡음 밖이 $3$ 개}다. 그리고 그
$3$ 개에 \textbf{이 사이클이 두 노트를 들여 설명한 것이 안 들어 있다.}

\begin{center}
\begin{tabular}{lrl}
\hline
주장 & 값 & $95\%$ 구간 \\
\hline
애니 트렌드 묶음이 약해진다(노트 492) & $-0.0592$ & $[-0.1052, +0.0150]$ \\
늘 붙던 행은 안 늙었다(노트 493) & $-0.0141$ & $[-0.1152, +0.0926]$ \\
새로 찬 행만 늙었다(노트 493) & $-0.0634$ & $[-0.1145, +0.0378]$ \\
아이돌이 제일 크다(노트 494) & $-0.3275$ & $[-0.4410, +0.1339]$ \\
\hline
\end{tabular}
\end{center}

\textbf{넷 다 $0$ 을 문다. 취소한다.}
\end{abstract}

\begin{figure}[t]
\centering
\includegraphics[width=\textwidth]{figs/f1.pdf}
\caption{\textbf{가} --- 일곱 도메인. \textbf{나} --- 노트 493 의 대조.
\textbf{다} --- 두 자의 잡음 폭.}
\end{figure}

\section{왜 안 보였나}

설명하려던 것은 \textbf{짝지은} 측정이었다 --- 노트 490 의 애니
$+0.0328$ 은 같은 씨앗 $12$ 개로 챔피언과 변형을 나란히 돌려 얻은
것이고 부호가 $12/12$ 다. 설명에 쓴 것은 \textbf{짝 없는} 측정이었다 ---
학습 행에서 잰 상관과 유보 행에서 잰 상관의 차이고, 두 표본이 다르므로
잡음이 통째로 남는다.

\textbf{그 자의 구간 폭이 $0.120$ 이다.} 설명하려던 효과($0.033$)의
\textbf{세 배가 넘는다.} 애초에 볼 수 없는 자로 보려 했다.

\begin{quote}
\textbf{조항} --- \textbf{짝지은 설계에서 나온 효과는 짝 없는 자로
설명하지 않는다.} 설명을 시작하기 전에 \emph{설명하려는 값이 그 자의
구간보다 큰지} 먼저 잰다. 노트 476 의 사촌이다 --- 거기서는 ``예보를 거의
안 바꾼다''를 무해함의 근거로 못 쓴다고 적었고, 여기서는 \emph{짝 없는
상관}을 설명의 근거로 못 쓴다고 적는다.
\end{quote}

\section{그래도 남는 것}

\textbf{셈은 잡음이 없다.} 다음 셋은 추정이 아니라 세어 나온 값이라 그대로
선다.

\begin{itemize}
\item \textbf{애니 트렌드 덮음이 학습 $56.4\%$ 에서 유보 $96.5\%$ 로
뛴다} --- 데이터랩 소급 한계가 아니다(애니 $2016$ 이전 $0$ 행).
\item \textbf{덮음이 $15\%$p 넘게 움직이는 짝이 $25$ 개다}(노트 494) ---
가드 \texttt{\_coverage\_shift} 는 그대로 둔다. 애초에 방향을 주장하지
않도록 썼다.
\item \textbf{노트 490 $\cdot$ 491 은 안 흔들린다} --- 시계열 수집 판
$-0.0283$($12/12$ 짝 씨앗)과 절단 절벽(정확히 $20$ 자 $18.3\%$ 대 $21$ 자
$0.6\%$, $p = 1.1 \times 10^{-22}$)은 각각 짝 설계와 큰 표본이다.
\end{itemize}

그리고 \textbf{애니가 시계열을 빼면 좋아진다는 사실 자체도 남는다}
($+0.0328$, $12/12$). 취소한 것은 \emph{왜}이지 \emph{무엇}이 아니다.

\section{잡음 밖인 셋}

\begin{center}
\begin{tabular}{llrr}
\hline
도메인 & 축 & 차 & 구간 \\
\hline
만화 & \texttt{goods\_scale} & $-0.2825$ & $[-0.3914, -0.0325]$ \\
세계애니 & \texttt{trend\_momentum} & $+0.1910$ & $[+0.0361, +0.3149]$ \\
아이돌 & \texttt{venue\_prominence} & $+0.2731$ & $[+0.0293, +0.4980]$ \\
\hline
\end{tabular}
\end{center}

묶음으로 보면 \textbf{세계애니 트렌드}($+0.1353$, $[+0.0200, +0.2481]$)
하나가 일곱 도메인 중 유일하게 $0$ 을 안 문다 --- \emph{강해지는} 쪽이다.

\section*{남은 것}

\begin{itemize}
\item \textbf{애니 $+0.0328$ 은 아직 설명이 없다} --- 세 후보(부호 ·
계절 · 늙음)를 다 떨어뜨렸고 넷째는 잡음이었다. \textbf{다음 시도는 짝
설계여야 한다} --- 예컨대 애니 행만 트렌드를 가린 모형과 챔피언을 같은
씨앗으로.
\item \textbf{만화 \texttt{goods\_scale} $-0.2825$} --- 잡음 밖인 셋 중
유일한 음수. 만화는 판 무게가 작지 않다.
\item \textbf{노트 494 의 문턱을 $10\%$p 로 낮추면} 짝이 늘지만 잡음
밖은 안 늘 것이다. \textbf{세는 대신 재야 한다.}
\end{itemize}

\end{document}
